\section{Related Work}
\label{sec:related}

\para{Phonetic representations in language models.}
The question of whether text-only models learn phonetic information has attracted growing interest.
\citet{mclaughlin2025rhyme} provided the most rigorous mechanistic study to date, identifying a ``phoneme mover head'' in \llama that promotes phonetically similar tokens during rhyme generation, with linear probes achieving approximately 96\% accuracy on embedding-layer phoneme prediction.
Their work demonstrated that phonetic information is linearly accessible in the embedding layer, but did not systematically track how this information evolves across all layers of the residual stream.
We extend their analysis by probing every layer of \gptmed, revealing that phonetic information is front-loaded rather than progressively built.
\citet{bunzeck2024small} compared grapheme-based and phoneme-based character-level language models on phonetic tasks, finding that even small models achieve 88--92\% rhyme prediction accuracy, with grapheme models surprisingly outperforming phoneme models.
Earlier work on character embeddings \citep{elhage2022interpreting} established that sub-word representations carry perceptual properties including sound-related features.

\para{Phonetic word embeddings.}
\citet{zouhar2023pwesuite} developed a standardized evaluation suite for phonetic word embeddings, including rhyme detection as a benchmark task.
Their work uses count-based, autoencoder, and contrastive learning methods with articulatory features from PanPhon, providing complementary evaluation methodology for measuring phonetic content in learned representations.
Our direction-vector analysis relates to this line of work: rather than training specialized phonetic embeddings, we analyze the implicit phonetic structure in general-purpose language model representations.

\para{Rhyme in neural text generation.}
Several systems have addressed rhyme explicitly in the context of poetry or song generation.
\citet{lau2018deep} proposed a joint neural model of poetic language, meter, and rhyme for sonnet generation.
\citet{jhamtani2019learning} used structured adversarial training to enforce rhyming constraints without manual phonetic features.
\citet{xue2021deeprapper} built a transformer-based rap generation system with explicit rhyme and rhythm modeling.
\citet{haider2018supervised} developed Siamese recurrent networks for supervised rhyme detection.
These systems treat rhyme as an explicit constraint to be enforced; our work instead asks where rhyme knowledge \emph{already exists} within standard pretrained models.

\para{Mechanistic interpretability of linguistic features.}
Our probing methodology follows the broader tradition of using linear probes to identify linguistic features in neural network representations \citep{radford2019language}.
While most probing studies find that syntactic and semantic features build through the network---with syntax peaking in middle layers and semantics in later layers---we find a qualitatively different pattern for phonetic features: they are strongest at the embedding boundary and gradually degrade.
This contributes to a more complete picture of how different levels of linguistic structure are distributed across transformer layers.
